% Luc Destombe --- licence CC BY-NC-SA 4.0
% https://creativecommons.org/licenses/by-nc-sa/4.0/deed.fr
% Fichier autonome : compiler avec pdflatex, deux fois (signets, nombre de pages).
\def\lyceeversion{1.0}
\documentclass[11pt,a4paper]{article}
\makeatletter
\newif\ifeleve
\elevefalse
\elevetrue

\newif\iflycee@palatino
\lycee@palatinotrue

\RequirePackage[utf8]{inputenc}
\RequirePackage[T1]{fontenc}
\RequirePackage[french]{babel}
\RequirePackage{etoolbox}

\newcommand{\ShorthandsOff}{\@ifundefined{shorthandoff}{}{\shorthandoff{;:!?}}}
\newcommand{\ShorthandsOn}{\@ifundefined{shorthandon}{}{\shorthandon{;:!?}}}
\ShorthandsOff
\AtBeginDocument{\ShorthandsOn}
\AtBeginEnvironment{tikzpicture}{\ShorthandsOff}

\RequirePackage{geometry}
\RequirePackage{fancyhdr}
\RequirePackage{lastpage}
\RequirePackage{multicol}
\RequirePackage{array}
\RequirePackage{enumitem}
\RequirePackage{xcolor}
\RequirePackage{needspace}

\RequirePackage{amsmath,amssymb}
\RequirePackage{mathpazo}
\renewcommand{\ttdefault}{lmtt}
\DeclareMathSizes{10.95}{10.95}{8}{5.5}
\begingroup\catcode`\_=\active \catcode`\^=\active
\gdef_#1{\sb{\scriptscriptstyle #1}}
\gdef^#1{\sp{\scriptscriptstyle\lycee@moinsepais #1}}
\endgroup
\newcommand\lycee@moins{\mathbin{%
  \setbox\z@\hbox{$\m@th\scriptscriptstyle\lycee@moinsorig$}%
  \dimen@=.55\wd\z@ \advance\dimen@-.5pt
  \hbox to.75\wd\z@{\hss$\m@th\scriptscriptstyle\vcenter{\hbox{%
    \tikz\draw[line width=.5pt,line cap=round](0,0)--(\the\dimen@,0);}}$\hss}}}
\begingroup\lccode`\~=`\- \lowercase{\endgroup\let~}\lycee@moins
\newcommand\lycee@moinsepais{\mathcode`\-="8000 }
\AtBeginDocument{\mathchardef\lycee@moinsorig=\mathcode`\-\relax}
\AtBeginDocument{\catcode`\_=12 \mathcode`\_="8000
  \catcode`\^=12 \mathcode`\^="8000 }
\newcommand\lycee@exposants{%
  \fontdimen13\textfont2=.48\fontdimen6\textfont2
  \fontdimen14\textfont2=.48\fontdimen6\textfont2
  \fontdimen15\textfont2=.48\fontdimen6\textfont2
  \fontdimen13\scriptfont2=.48\fontdimen6\scriptfont2
  \fontdimen14\scriptfont2=.48\fontdimen6\scriptfont2
  \fontdimen15\scriptfont2=.48\fontdimen6\scriptfont2}
\every@math@size\expandafter{\the\every@math@size\lycee@exposants}
\newlength{\retraitcalculs}
\setlength{\retraitcalculs}{2em}

\RequirePackage{tikz}
\usetikzlibrary{arrows.meta,calc,babel,shapes.misc}
\RequirePackage[most]{tcolorbox}
\tcbuselibrary{skins,breakable}

\definecolor{coulDef}{HTML}{1F4E79}
\definecolor{coulProp}{HTML}{7030A0}
\definecolor{coulRem}{HTML}{C55A11}
\definecolor{coulEx}{HTML}{2E7D32}
\definecolor{coulExo}{HTML}{B03A2E}
\definecolor{coulPart}{HTML}{1F4E79}
\definecolor{coulMeth}{HTML}{1B6E4F}
\definecolor{coulCourbe}{HTML}{0B5ED7}
\definecolor{coulCourbeB}{HTML}{C00000}
\definecolor{coulCourbeC}{HTML}{2E7D32}
\definecolor{coulGrille}{HTML}{8C8C8C}
\definecolor{coulRep}{HTML}{1B6E4F}
\definecolor{coulBar}{HTML}{2E7D32}

\newcommand{\R}{\mathbb{R}}
\newcommand{\Rs}{\mathbb{R}^{*}}
\renewcommand{\le}{\leqslant}
\renewcommand{\ge}{\geqslant}
\renewcommand{\approx}{\simeq}
\newcommand{\vect}[1]{\overrightarrow{#1}}
\newcommand{\norme}[1]{\left\lVert #1 \right\rVert}
\newcommand{\intervalle}[2]{\left[#1\,;#2\right]}
\RequirePackage{cancel}
\renewcommand{\CancelColor}{\color{coulExo}}
\newcommand{\fois}[1]{\times{\color{coulExo}#1}}
\newcommand{\barre}[1]{\cancel{#1}}
\newlength{\lycee@eqh}\newlength{\lycee@eqd}
\newsavebox{\lycee@eqa}\newsavebox{\lycee@eqb}\newsavebox{\lycee@eqc}
\newcommand{\lycee@eqcol}[1]{\hbox{$\renewcommand{\arraystretch}{1.35}%
  \begin{array}[t]{@{}>{\displaystyle}r@{}>{\displaystyle{}}l@{}}#1\end{array}$}}
\newcommand{\lycee@eqmesure}[1]{%
  \ifdim\ht#1>\lycee@eqh\setlength{\lycee@eqh}{\ht#1}\fi
  \ifdim\dp#1>\lycee@eqd\setlength{\lycee@eqd}{\dp#1}\fi}
\newcommand{\lycee@eqtrait}{\hspace{1.2em}\tikz[baseline=0pt]\draw[dash pattern=on 2pt off 2pt,
  line width=0.5pt] (0,\lycee@eqh) -- (0,-\lycee@eqd);\hspace{1.2em}}
\newcommand{\eqproduit}[2]{\par\smallskip\noindent C'est une équation produit :
  $A\times B=0 \iff A=0$ ou $B=0$.\par\nopagebreak\smallskip
  \begingroup\setlength{\lycee@eqh}{0pt}\setlength{\lycee@eqd}{0pt}%
  \setbox\lycee@eqa=\lycee@eqcol{#1}\setbox\lycee@eqb=\lycee@eqcol{#2}%
  \lycee@eqmesure\lycee@eqa\lycee@eqmesure\lycee@eqb%
  \noindent\hspace*{\retraitcalculs}\usebox\lycee@eqa\lycee@eqtrait\usebox\lycee@eqb\endgroup\par\smallskip}
\newcommand{\eqproduittrois}[3]{\par\smallskip\noindent C'est une équation produit :
  $A\times B\times C=0 \iff A=0$ ou $B=0$ ou $C=0$.\par\nopagebreak\smallskip
  \begingroup\setlength{\lycee@eqh}{0pt}\setlength{\lycee@eqd}{0pt}%
  \setbox\lycee@eqa=\lycee@eqcol{#1}\setbox\lycee@eqb=\lycee@eqcol{#2}\setbox\lycee@eqc=\lycee@eqcol{#3}%
  \lycee@eqmesure\lycee@eqa\lycee@eqmesure\lycee@eqb\lycee@eqmesure\lycee@eqc%
  \noindent\hspace*{\retraitcalculs}\usebox\lycee@eqa\lycee@eqtrait\usebox\lycee@eqb\lycee@eqtrait\usebox\lycee@eqc\endgroup
  \par\smallskip}

\tikzset{grille/.style={coulGrille,line width=0.4pt}}
\newcommand{\quadrillage}[4]{\draw[grille,step=1] (#1,#2) grid (#3,#4);}
\tikzset{
  croix/.style={cross out,draw,line width=0.6pt,inner sep=0pt,minimum size=4pt},
  croixdroite/.style={croix,rotate=45,line width=0.8pt},
}
\newcommand{\pt}[3]{\path (#1) node[croix]{} node[#2,font=\small,inner sep=2.5pt]{$#3$};}

\newcommand{\lycee@repere}[4]{%
  \draw[grille,step=1] (#1,#2) grid (#3,#4);
  \draw[-{Stealth[length=2mm]},thick] (#1,0)--({#3+0.4},0) node[below left=-1pt]{$x$};
  \draw[-{Stealth[length=2mm]},thick] (0,#2)--(0,{#4+0.4}) node[below left=-1pt]{$y$};
  \node[below left,font=\scriptsize,inner sep=1.5pt] at (0,0) {$0$};}
\newcommand{\lycee@gradx}[1]{\foreach \k/\l in {#1}{%
  \draw (\k,0.07)--(\k,-0.07) node[below,font=\scriptsize,inner sep=1.5pt]{$\l$};}}
\newcommand{\lycee@grady}[1]{\foreach \k/\l in {#1}{%
  \draw (0.07,\k)--(-0.07,\k) node[left,font=\scriptsize,inner sep=1.5pt]{$\l$};}}
\newcommand{\lycee@intff}[2]{\left[#1\,;#2\right]}
\newcommand{\lycee@intfo}[2]{\left[#1\,;#2\right[}
\newcommand{\lycee@intof}[2]{\left]#1\,;#2\right]}
\newcommand{\lycee@intoo}[2]{\left]#1\,;#2\right[}
\newcommand{\lycee@ens}[1]{\left\{#1\right\}}
\AtBeginDocument{%
  \providecommand{\repere}{\lycee@repere}%
  \providecommand{\gradx}{\lycee@gradx}%
  \providecommand{\grady}{\lycee@grady}%
  \providecommand{\Cf}{\mathcal{C}\sb{\scriptscriptstyle f}}%
  \providecommand{\Cg}{\mathcal{C}\sb{\scriptscriptstyle g}}%
  \providecommand{\intff}{\lycee@intff}%
  \providecommand{\intfo}{\lycee@intfo}%
  \providecommand{\intof}{\lycee@intof}%
  \providecommand{\intoo}{\lycee@intoo}%
  \providecommand{\ens}{\lycee@ens}}

\newlist{questions}{enumerate}{3}
\setlist[questions,1]{label=\arabic*\textdegree),leftmargin=*,itemsep=2pt,topsep=3pt}
\setlist[questions,2]{label=\alph*),leftmargin=*,itemsep=1pt,topsep=2pt}
\setlist[questions,3]{label=\roman*),leftmargin=*,itemsep=1pt,topsep=1pt}
\setlist[itemize]{label=\textbullet,leftmargin=*,itemsep=2pt,topsep=3pt}

\newcommand{\besoinplace}[1]{\par\penalty\@M\begingroup
  \dimen@=#1\relax\dimen@ii=\pagegoal\advance\dimen@ii-\pagetotal
  \ifdim\dimen@>\dimen@ii\ifdim\dimen@ii>\z@\vfil\fi\break\fi\endgroup}

\newcommand{\lycee@pied}{\thepage/\pageref{LastPage}}
\newcommand{\entete}[2]{%
  \pagestyle{fancy}\fancyhf{}%
  \fancyhead[L]{\footnotesize\color{coulPart}\textbf{#1}}%
  \fancyhead[R]{\footnotesize\color{coulPart}\textbf{#2}}%
  \fancyfoot[C]{\footnotesize\lycee@pied}%
  \renewcommand{\headrulewidth}{0.4pt}%
  \renewcommand{\headrule}{\hbox to\headwidth{%
    \color{coulPart!40}\leaders\hrule height \headrulewidth\hfill}}}

\RequirePackage{ccicons}
\newcommand{\lycee@licence}{{\scriptsize\color{gray}%
  \copyright\ Luc Destombe --- \ccbyncsa\ CC BY-NC-SA 4.0}}
\newif\iflycee@licence \lycee@licencetrue
\newcommand{\sanslicence}{\lycee@licencefalse}
\AtBeginDocument{\iflycee@licence\AddToHookNext{shipout/foreground}{%
  \put(0,0){\rlap{\kern\dimexpr1in+\hoffset+\oddsidemargin+\textwidth\relax
    \llap{\raisebox{-\dimexpr1in+\voffset+\topmargin+\headheight+\headsep
      +\textheight+\footskip\relax}{\lycee@licence}}}}}\fi}

\newcommand{\lycee@site}{faire-des-maths.fr}
\newcommand{\lycee@version}{{\scriptsize\color{gray}\lycee@site
  \ifdefined\lyceeversion\ --- Version~\lyceeversion\fi}}
\AtBeginDocument{\iflycee@licence\AddToHookNext{shipout/foreground}{%
  \put(0,0){\rlap{\kern\dimexpr1in+\hoffset+\oddsidemargin\relax
    \raisebox{-\dimexpr1in+\voffset+\topmargin+\headheight+\headsep
      +\textheight+\footskip\relax}{\lycee@version}}}}\fi}

\newcommand{\formulesserrees}{%
  \setlength{\abovedisplayskip}{3pt}\setlength{\belowdisplayskip}{3pt}%
  \setlength{\abovedisplayshortskip}{2pt}\setlength{\belowdisplayshortskip}{3pt}}

\setlength{\parindent}{0pt}

\RequirePackage[implicit=false,hidelinks,pdfencoding=auto,psdextra,bookmarksopen,
  bookmarksopenlevel=1,pdfcreator={},pdfproducer={}]{hyperref}
\RequirePackage{bookmark}
\hypersetup{pdfauthor={Luc Destombe},pdfkeywords={CC BY-NC-SA 4.0}}
\newcounter{lycee@signet}
\newcommand{\signet}[3][]{\stepcounter{lycee@signet}%
  \edef\lycee@dest{\if\relax\detokenize{#1}\relax
    signet.\arabic{lycee@signet}\else#1\fi}%
  \hypertarget{\lycee@dest}{}\bookmark[level=#2,dest=\lycee@dest]{#3}}
\pdfstringdefDisableCommands{%
  \def\R{R}\def\vect#1{#1}\def\Cf{Cf}\def\Cg{Cg}\def\fois#1{×#1}%
  \def\barre#1{#1}\def\intervalle#1#2{[#1;#2]}\def\intff#1#2{[#1;#2]}%
  \def\textsuperscript#1{#1}\def\dfrac#1#2{#1/#2}\def\frac#1#2{#1/#2}%
  \def\vec#1{#1⃗}}%
\newcommand{\res}[1]{#1}
\geometry{top=2cm,bottom=1cm,left=1cm,right=1cm,headheight=14pt,footskip=0.6cm}

\RequirePackage{pgfplots}
\pgfplotsset{compat=1.18}
\RequirePackage{tkz-tab}

\newcounter{partiecours}
\renewcommand{\thepartiecours}{\Roman{partiecours}}
\newcommand{\partiecours}[1]{%
  \refstepcounter{partiecours}%
  \Needspace*{8\baselineskip}%
  \bigskip
  \noindent\begin{tcolorbox}[enhanced,colback=coulPart,colframe=coulPart,
      boxrule=0pt,arc=2pt,left=8pt,right=8pt,top=4pt,bottom=4pt,
      before skip=6pt,after skip=10pt]
    \color{white}\bfseries\large \thepartiecours{} --- #1%
    \signet[partie-\arabic{partiecours}]{1}{\thepartiecours{} --- #1}
  \end{tcolorbox}%
  \addcontentsline{toc}{section}{\thepartiecours{} --- #1}%
}

\newtcolorbox{boitecours}[3][]{%
  enhanced,breakable,
  colback=white,colframe=#2,coltitle=white,
  fonttitle=\bfseries,
  attach boxed title to top left={xshift=8pt,yshift=-\tcboxedtitleheight/2},
  boxed title style={colback=#2,arc=2pt,boxrule=0pt},
  boxrule=0.9pt,arc=2pt,
  left=8pt,right=8pt,top=8pt,bottom=6pt,
  title={#3},#1}

\newcounter{methode}
\newenvironment{definitionbox}[1][Définition]%
  {\begin{boitecours}[phantom={\signet{2}{#1}}]{coulDef}{#1}}{\end{boitecours}}
\newenvironment{proprietebox}[1][Propriété]%
  {\begin{boitecours}[phantom={\signet{2}{#1}}]{coulProp}{#1}}{\end{boitecours}}
\newenvironment{methodebox}[1][Méthode]{\stepcounter{methode}%
  \begin{boitecours}[phantom={\signet[methode-\arabic{methode}]{2}{#1}}]{coulMeth}{#1}}%
  {\end{boitecours}}

\newcommand{\etape}[1]{\par\smallskip\textbf{\color{coulMeth}#1}\par\nobreak\smallskip}
\newcommand{\enonce}[1]{\emph{#1}\par\smallskip}
\newcommand{\reponse}[1]{\par\smallskip
  {\footnotesize\itshape\color{black!55}Réponses : #1}}
\newcommand{\demo}[1]{\textbf{\color{coulMeth}Démonstration.} #1}
\newcommand{\afaire}[1]{%
  \par\smallskip
  \noindent{\small\color{coulExo}$\blacktriangleright$\ \textbf{À faire :}
    fiche d'exercices du chapitre, #1.}\par\smallskip}

\newenvironment{calculs}{\setlength{\jot}{5pt}\@fleqntrue\@mathmargin\retraitcalculs
  \setlength{\abovedisplayskip}{4pt}\setlength{\belowdisplayskip}{4pt}%
  \setlength{\abovedisplayshortskip}{4pt}\setlength{\belowdisplayshortskip}{4pt}%
  \csname alignat*\endcsname{2}}%
  {\csname endalignat*\endcsname}
\newcommand{\rem}[1]{\text{\small\itshape\color{black!60}#1}}
\newcommand{\explique}[2]{\par\smallskip\noindent
  \begin{minipage}[t]{0.57\linewidth}\raggedright #1\end{minipage}\hfill
  \begin{minipage}[t]{0.40\linewidth}\raggedright\leavevmode
    {\small\itshape\color{black!60}#2}\end{minipage}\par}
\newcommand{\cas}[1]{\par\smallskip\noindent\textbf{\color{coulExo}#1}\vadjust{\nobreak}}
\newcommand{\calc}[1]{\par\smallskip\textbf{\color{coulExo}#1}\par\nobreak}

\newtcolorbox{remarquebox}[1][Remarque]{%
  enhanced,breakable,colback=white,colframe=coulRem,
  boxrule=0pt,leftrule=2.5pt,arc=0pt,outer arc=0pt,
  left=8pt,right=6pt,top=5pt,bottom=5pt,
  before upper={\textbf{\color{coulRem}#1}\par\smallskip}}

\newtcolorbox{exemplebox}[1][Exemple]{%
  enhanced,breakable,colback=white,colframe=coulEx,
  boxrule=0pt,leftrule=2.5pt,arc=0pt,outer arc=0pt,
  left=8pt,right=6pt,top=5pt,bottom=5pt,
  before upper={\textbf{\color{coulEx}#1}\par\smallskip}}

\pgfplotsset{
  repere/.style={
    axis lines=middle,
    axis line style={-{Stealth[length=2.4mm]},thick},
    every axis x label/.style={at={(ticklabel* cs:1.0)},anchor=west},
    every axis y label/.style={at={(ticklabel* cs:1.0)},anchor=south},
    label style={font=\footnotesize},
    tick label style={font=\scriptsize},
    grid=both,
    major grid style={grille},
    minor tick num=0,
    tick align=inside,
    clip mode=individual,
  },
}
\tikzset{
  courbe/.style={coulCourbe,line width=1.1pt,smooth},
  courbeB/.style={coulCourbeB,line width=1.1pt,smooth},
  courbeC/.style={coulCourbeC,line width=1.1pt,smooth},
}

\newlist{etapes}{enumerate}{1}
\setlist[etapes]{label=\textbf{\arabic*.},leftmargin=*,itemsep=1pt,topsep=2pt}

\newlist{expressions}{enumerate}{1}
\setlist[expressions]{label={\Alph*\ $=$},leftmargin=*,itemsep=3pt,topsep=3pt}

\newcounter{exo}
\renewcommand{\theexo}{\ifnum\value{exo}<10 0\fi\arabic{exo}}
\newtcolorbox{exercice}[1][]{%
  enhanced,breakable,
  colback=white,colframe=coulExo,
  boxrule=0.9pt,arc=2pt,
  left=8pt,right=8pt,top=8pt,bottom=6pt,
  coltitle=white,fonttitle=\bfseries,
  attach boxed title to top left={xshift=8pt,yshift=-\tcboxedtitleheight/2},
  boxed title style={colback=coulExo,arc=2pt,boxrule=0pt},
  phantom={\signet{2}{Exercice \arabic{exo}}},
  title={Exercice \theexo\ifx\relax#1\relax\else{} --- #1\fi}}
\newenvironment{exo}[1][\relax]{\refstepcounter{exo}\begin{exercice}[#1]}{\end{exercice}}

  \RequirePackage{comment}
  \excludecomment{correction}
  \renewcommand{\reponse}[1]{}
  \newcommand{\autableau}{\hfill{\footnotesize\itshape\color{black!45}(au tableau)}}
  \newcommand{\etapeexemples}{%
    \par\smallskip\textbf{\color{coulMeth}Exemples}\autableau\par\nobreak\smallskip}
  \newcommand{\etapetableau}[1]{%
    \par\smallskip\textbf{\color{coulMeth}#1}\autableau\par\nobreak\smallskip}
  \newtcolorbox{exempletableau}[1][Exemple]{%
    enhanced,breakable,colback=white,colframe=coulEx,
    boxrule=0pt,leftrule=2.5pt,arc=0pt,outer arc=0pt,
    left=8pt,right=6pt,top=4pt,bottom=4pt,
    before skip=5pt,after skip=5pt,
    before upper={\textbf{\color{coulEx}#1}\autableau\par\smallskip}}

\RequirePackage{comment}
\excludecomment{correctionavj}

\newcommand{\coursEntete}{}
\newcommand{\coursNiveau}{}
\newcommand{\coursTitre}{}
\newcommand{\coursSurTitre}{}
\newcommand{\coursSousTitre}{}

\newcommand{\enteteCours}[1]{\renewcommand{\coursEntete}{#1}}
\newcommand{\chapitrecours}[2]{\enteteCours{Chapitre #1 --- #2}}
\newcommand{\niveaucours}[1]{\renewcommand{\coursNiveau}{#1}}
\newcommand{\titrecours}[3][]{%
  \renewcommand{\coursTitre}{#2}%
  \renewcommand{\coursSurTitre}{#1}%
  \renewcommand{\coursSousTitre}{#3}}

\pagestyle{fancy}
\fancyhf{}
\fancyhead[L]{\footnotesize\color{coulPart}\textbf{\coursEntete}}
\fancyhead[R]{\footnotesize\color{coulPart}\textbf{\coursNiveau}}
\fancyfoot[C]{\footnotesize\thepage}
\renewcommand{\headrulewidth}{0.4pt}
\renewcommand{\headrule}{\hbox to\headwidth{\color{coulPart!40}\leaders\hrule height \headrulewidth\hfill}}

\setlength{\parskip}{2pt}

\newcommand{\lycee@mentioneleve}{\\[2pt]{\small\itshape Version élève}}
\AtBeginDocument{%
  \begin{center}
    \begin{tcolorbox}[enhanced,width=0.72\linewidth,colback=white,colframe=coulPart,
        boxrule=1.6pt,arc=3pt,halign=center,top=10pt,bottom=10pt]
      \ifx\coursSurTitre\empty
        {\Huge\bfseries\color{coulPart} \coursTitre}\\[4pt]
      \else
        {\Huge\bfseries\color{coulPart} \coursTitre}\\[3pt]
        {\Large\bfseries\color{coulPart} \coursSurTitre}\\[5pt]
      \fi
      {\normalsize \coursSousTitre}%
      \lycee@mentioneleve
    \end{tcolorbox}
  \end{center}%
}

\RequirePackage{tikz}
\RequirePackage{xcolor}

\usetikzlibrary{decorations.text}
\definecolor{lycee@metal}{HTML}{1F4E79}
\definecolor{lycee@manche}{HTML}{8A5A2B}
\definecolor{lycee@etincelle}{HTML}{F28C28}
\newcommand{\lycee@pluriel}[2]{#1 #2\ifnum#1>1 s\fi}
\newcommand{\lycee@sceau}{%
  \edef\lycee@bas{\lycee@pluriel{\lyceerelectures}{relecture} et
    \lycee@pluriel{\lyceerodages}{rodage}}%
  \begin{tikzpicture}[line join=round]
    \fill[coulPart!10,draw=coulPart,line width=.8pt] (0:2.2)
      \foreach \i in {1,...,71}{ -- (\i*5:{isodd(\i)?2.1:2.2}) } -- cycle;
    \draw[coulPart,line width=1.2pt] (0,0) circle (1.93);
    \draw[coulPart,line width=.4pt] (0,0) circle (1.85);
    \draw[coulPart,line width=.4pt] (0,0) circle (1.2);
    \path[decorate,decoration={text along path,text align=center,
        text={|\fontsize{9.5}{11}\selectfont\bfseries\color{coulPart}|TEST{\'E} AVEC MES {\'E}L{\`E}VES}}]
      (180:1.45) arc (180:0:1.45);
    \path[decorate,decoration={text along path,text align=center,
        text/.expanded={|\noexpand\fontsize{8.5}{10}\noexpand\selectfont\noexpand\color{coulPart}|\lycee@bas}}]
      (195:1.68) arc (195:345:1.68);
    \begin{scope}[x=.25mm,y=-.25mm]
      \begin{scope}[shift={(-32,-40)}]%
        \fill[lycee@metal] (23,31) rectangle (41,49);
        \draw[lycee@etincelle,line width=.5pt,line cap=round]
          (20,30)--(15,27) (19,34)--(14,35) (44,30)--(49,27) (45,34)--(50,35);
        \fill[lycee@manche,rounded corners=.3mm] (38,24.6) rectangle (60,27.6);
        \fill[lycee@metal,rounded corners=.4mm] (25,20) rectangle (39,31);
        \fill[lycee@etincelle,rounded corners=.2mm] (25,20) rectangle (28,31);
      \end{scope}
    \end{scope}
    \ifdefined\lyceesceauplusieurs\def\lycee@annee{Ann{\'e}es~:}\else
      \def\lycee@annee{Ann{\'e}e}\fi
    \path[decorate,decoration={text along path,text align=center,
        text/.expanded={|\noexpand\fontsize{8}{10}\noexpand\selectfont\noexpand\color{coulPart}|%
          \unexpanded\expandafter{\lycee@annee} \unexpanded\expandafter{\lyceesceauannees}}}]
      (95:1.08) arc (95:445:1.08);
  \end{tikzpicture}}
\AtEndDocument{\ifdefined\lyceesceauannees
  \par\besoinplace{5cm}\vspace{1.5em}%
  \begin{center}\lycee@sceau\end{center}\fi}
\makeatother

\chapitrecours{13}{Systèmes d'équations}
\niveaucours{Seconde}
\titrecours{SYSTÈMES D'ÉQUATIONS}{Cours --- Seconde}

\begin{document}

\ShorthandsOff
\tikzset{
  droite/.style={coulCourbe,line width=1.1pt},
  droiteB/.style={coulCourbeB,line width=1.1pt},
  aide/.style={dashed,black!45,line width=0.6pt},
}
\ShorthandsOn
\newcommand{\sys}[1]{\left\{\begin{aligned}#1\end{aligned}\right.}
\newcommand{\coord}[2]{\begin{pmatrix}#1\\#2\end{pmatrix}}

\partiecours{Systèmes par substitution}

Au cinéma, Lina paie $31$~€ pour deux places adultes et une place enfant ; Tom paie
$19$~€ pour une place adulte et une place enfant. On note $x$ le prix d'une place
adulte et $y$ celui d'une place enfant, en euros. Les deux phrases se traduisent par
deux équations, qui doivent être vraies \textbf{en même temps} :
\begin{center}
  ${\sys{2x+y &= 31 && (L_{1})\\ x+y &= 19 && (L_{2})}}$
\end{center}

\begin{definitionbox}[Définition 1 --- Système de deux équations à deux inconnues]
\begin{itemize}
  \item Deux équations d'inconnues $x$ et $y$, à vérifier en même temps, forment un
        \textbf{système}. On les nomme $(L_{1})$ et $(L_{2})$ (première et deuxième
        ligne).
  \item Une \textbf{solution} du système est un couple $(x\,;y)$ qui vérifie
        \textbf{les deux} équations.
  \item \textbf{Résoudre} le système, c'est trouver tous ses couples solutions ;
        on note $S$ leur ensemble.
\end{itemize}
\end{definitionbox}

\begin{exemplebox}
Pour le système du cinéma :
\begin{itemize}
  \item $(12\,;7)$ est une solution : $2\times 12+7=31$ et $12+7=19$ ;
  \item $(10\,;11)$ n'est pas une solution : $2\times 10+11=31$, mais $10+11=21$, et
        non $19$. Un couple qui ne vérifie qu'une seule équation ne convient pas.
\end{itemize}
\end{exemplebox}

\begin{methodebox}[Méthode 1 --- Résoudre un système par substitution]
Quand une inconnue a pour coefficient $1$ ou $-1$, on l'isole dans son équation, puis on
la remplace par son expression dans l'autre équation : il reste une équation à une seule
inconnue, que l'on sait résoudre.
\etapeexemples
Résoudre chaque système.
\begin{center}
  a)~${\sys{y &= 2x-1 && (L_{1})\\ 3x+y &= 9 && (L_{2})}}$ \qquad
  b)~${\sys{x+3y &= 7 && (L_{1})\\ 2x-y &= 7 && (L_{2})}}$ \qquad
  c)~${\sys{4x-y &= 9 && (L_{1})\\ 3x+2y &= 4 && (L_{2})}}$
\end{center}

\etape{À vous de jouer}
Résoudre chaque système.
\begin{center}
  a)~${\sys{y &= 3x+2 && (L_{1})\\ 2x+y &= 12 && (L_{2})}}$ \qquad
  b)~${\sys{x+2y &= 8 && (L_{1})\\ 3x-y &= 3 && (L_{2})}}$ \qquad
  c)~${\sys{3x-y &= 5 && (L_{1})\\ 2x+3y &= 7 && (L_{2})}}$
\end{center}

\end{methodebox}

\afaire{exercices 1 à 4}

\partiecours{Systèmes par combinaison}

Dans le système ${\sys{3x+4y &= 2\\ 2x+3y &= 1}}$, aucune inconnue n'a pour
coefficient $1$ ou $-1$ : en isolant $x$ ou $y$, on ferait apparaître des fractions.
On fait plutôt \textbf{disparaître} une inconnue en additionnant les deux équations.

\begin{proprietebox}[Propriété 1 --- Combinaison de deux équations]
Si un couple $(x\,;y)$ est solution d'un système, il vérifie aussi :
\begin{itemize}
  \item l'équation obtenue en multipliant les deux membres de $(L_{1})$, ou de
        $(L_{2})$, par un même nombre non nul ;
  \item l'équation obtenue en additionnant membre à membre $(L_{1})$ et $(L_{2})$.
\end{itemize}
Quand les coefficients d'une inconnue sont \textbf{opposés} ($3$ et $-3$), cette
inconnue disparaît de la somme.
\end{proprietebox}

\begin{methodebox}[Méthode 2 --- Résoudre un système par combinaison]
On multiplie au besoin une équation, ou les deux, pour que les coefficients d'une
inconnue soient opposés ; on additionne : il reste une équation à une inconnue. On
remplace ensuite la valeur trouvée dans $(L_{1})$ pour calculer l'autre inconnue.
\etapeexemples
Résoudre chaque système.
\begin{center}
  a)~${\sys{3x+2y &= 12 && (L_{1})\\ 5x-2y &= 4 && (L_{2})}}$ \qquad
  b)~${\sys{2x+3y &= 12 && (L_{1})\\ 4x+5y &= 22 && (L_{2})}}$ \qquad
  c)~${\sys{3x+4y &= 2 && (L_{1})\\ 2x+3y &= 1 && (L_{2})}}$
\end{center}

\etape{À vous de jouer}
Résoudre chaque système.
\begin{center}
  a)~${\sys{4x+3y &= 18 && (L_{1})\\ 2x-3y &= 0 && (L_{2})}}$ \qquad
  b)~${\sys{3x+2y &= 7 && (L_{1})\\ 6x+5y &= 16 && (L_{2})}}$ \qquad
  c)~${\sys{2x+5y &= 1 && (L_{1})\\ 3x+4y &= 5 && (L_{2})}}$
\end{center}

\end{methodebox}

\begin{remarquebox}[Choisir sa méthode]
Une inconnue a pour coefficient $1$ ou $-1$ : la substitution est la plus rapide.
Sinon, la combinaison évite les fractions.
\end{remarquebox}

\afaire{exercices 5 à 8}

\partiecours{Problèmes et systèmes}

\begin{methodebox}[Méthode 3 --- Résoudre un problème à l'aide d'un système]
On nomme les deux inconnues, avec leur unité ; on traduit chaque information par une
équation ; on résout le système, puis on répond à la question par une phrase.
\etapeexemples
\begin{questions}[label=\alph*)]
  \item Sur un parking, il y a $40$ véhicules, des voitures et des motos, et l'on compte
        $130$ roues. Combien y a-t-il de voitures et de motos ?
  \item Au cinéma, trois places adultes et deux places enfants coûtent $49$~€ ; deux
        places adultes et trois places enfants coûtent $46$~€. Quel est le prix d'une
        place adulte et celui d'une place enfant ?
\end{questions}

\etape{À vous de jouer}
\begin{questions}[label=\alph*)]
  \item Une salle de concert a vendu $300$ billets : des places assises à $25$~€ et des
        places debout à $15$~€. La recette est de $5\,900$~€. Combien de places de
        chaque sorte a-t-elle vendues ?
  \item Dans une salle d'escalade, trois entrées adultes et deux entrées enfants coûtent
        $54$~€ ; deux entrées adultes et trois entrées enfants coûtent $51$~€. Quel est
        le prix d'une entrée adulte et celui d'une entrée enfant ?
\end{questions}

\end{methodebox}

\afaire{exercices 9 à 11}

\partiecours{Systèmes et intersection de droites}

Rappel (chapitre 9) : la droite d'équation cartésienne $ax+by+c=0$ a pour vecteur
directeur $\vec{u}\coord{-b}{a}$ ; deux droites sont sécantes quand le déterminant de
leurs vecteurs directeurs n'est pas nul. Une équation $ax+by=c$ d'un système est donc
l'équation d'une droite.

\begin{proprietebox}[Propriété 2 --- Point d'intersection et système]
\begin{minipage}[c]{0.62\linewidth}\raggedright
Soit deux droites sécantes $(d)$ : $ax+by+c=0$ et $(d')$ : $a'x+b'y+c'=0$. Les
coordonnées $(x\,;y)$ de leur point d'intersection $K$ vérifient les deux équations :
c'est l'unique solution du système
\[
  \sys{ax+by &= -c\\ a'x+b'y &= -c'}
\]
\end{minipage}\hfill
\begin{minipage}[c]{0.34\linewidth}\centering
\begin{tikzpicture}[x=0.4cm,y=0.4cm]
  \draw[droite] (0,0.5) -- (6,4);
  \draw[droiteB] (0,4.5) -- (6,0.3);
  \path (3.55,2.57) node[croix]{} node[below=2pt,font=\small]{$K$};
  \node[coulCourbe,font=\footnotesize,above left] at (6,4) {$(d)$};
  \node[coulCourbeB,font=\footnotesize,above] at (0.6,4.4) {$(d')$};
\end{tikzpicture}
\end{minipage}
\end{proprietebox}

\begin{methodebox}[Méthode 4 --- Calculer les coordonnées du point d'intersection de deux droites]
On vérifie que les droites sont sécantes (déterminant de leurs vecteurs directeurs non
nul), puis on résout le système formé par leurs deux équations.
\etapeexemples
Pour chaque couple de droites, démontrer que les droites sont sécantes, puis calculer
les coordonnées de leur point d'intersection $K$.
\begin{questions}[label=\alph*)]
  \item $(d_{1})$ : $x+3y-7=0$ et $(d_{2})$ : $2x-y-7=0$.
  \item $(d_{3})$ : $y=2x-1$ et $(d_{4})$ : $2x+3y-13=0$.
\end{questions}

\etape{À vous de jouer}
Pour chaque couple de droites, démontrer que les droites sont sécantes, puis calculer
les coordonnées de leur point d'intersection $K$.
\begin{questions}[label=\alph*)]
  \item $(d_{1})$ : $2x+y-8=0$ et $(d_{2})$ : $x-3y+3=0$.
  \item $(d_{3})$ : $y=-x+4$ et $(d_{4})$ : $3x-2y-7=0$.
\end{questions}

\end{methodebox}

\afaire{exercices 12 à 15}

\partiecours{Équation réduite d'une droite par un système}

Rappel (chapitre 9) : le point $A(x_{A}\,;y_{A})$ est sur la droite d'équation $y=mx+p$
quand ses coordonnées vérifient l'équation : $y_{A}=m\times x_{A}+p$. Pour une droite
qui passe par deux points connus, chaque point donne une équation d'inconnues $m$ et
$p$ : on obtient un système.

\begin{methodebox}[Méthode 5 --- Trouver l'équation réduite d'une droite par un système]
On écrit que chaque point vérifie $y=mx+p$, puis on résout le système d'inconnues $m$ et
$p$ par combinaison : en multipliant $(L_{1})$ par $-1$, les $p$ s'éliminent.
\etapeexemples
Pour chaque couple de points, trouver l'équation réduite de la droite $(AB)$.
\begin{questions}[label=\alph*)]
  \item $A(1\,;3)$ et $B(4\,;9)$.
  \item $A(-1\,;2)$ et $B(2\,;4)$.
\end{questions}

\etape{À vous de jouer}
Pour chaque couple de points, trouver l'équation réduite de la droite $(AB)$.
\begin{questions}[label=\alph*)]
  \item $A(1\,;1)$ et $B(3\,;7)$.
  \item $A(-2\,;1)$ et $B(1\,;0)$.
\end{questions}

\end{methodebox}

\afaire{exercices 16 à 18}

\partiecours{Aucune ou une infinité de solutions}

\begin{proprietebox}[Propriété 3 --- Nombre de solutions d'un système]
Les deux équations d'un système sont celles de deux droites $(d)$ et $(d')$.
\begin{center}
\begin{tabular}{ccc}
\begin{tikzpicture}[x=0.38cm,y=0.38cm]
  \draw[droite] (0,0.5) -- (5,3.5);
  \draw[droiteB] (0,3.5) -- (5,0.5);
  \path (2.5,2) node[croix]{};
\end{tikzpicture}
&
\begin{tikzpicture}[x=0.38cm,y=0.38cm]
  \draw[droite] (0,0.3) -- (5,2.3);
  \draw[droiteB] (0,1.7) -- (5,3.7);
\end{tikzpicture}
&
\begin{tikzpicture}[x=0.38cm,y=0.38cm]
  \draw[coulCourbe,line width=2.6pt] (0,0.8) -- (5,3.2);
  \draw[coulCourbeB,line width=1pt] (0,0.8) -- (5,3.2);
\end{tikzpicture}
\\[2pt]
sécantes : & strictement parallèles : & confondues : \\
\textbf{une seule solution} & \textbf{aucune solution} & \textbf{une infinité}\\
 & $S=\varnothing$ & \textbf{de solutions}
\end{tabular}
\end{center}
Pendant la résolution, si les deux inconnues disparaissent, il reste :
\begin{itemize}
  \item une égalité \textbf{fausse} (comme $-6=1$) : le système n'a \textbf{aucune
        solution} ;
  \item une égalité \textbf{toujours vraie} (comme $12=12$) : le système a une
        \textbf{infinité de solutions}, tous les couples $(x\,;y)$ qui vérifient
        $(L_{1})$.
\end{itemize}
\end{proprietebox}

\begin{methodebox}[Méthode 6 --- Reconnaître un système sans solution ou avec une infinité de solutions]
On résout comme d'habitude ; si les inconnues disparaissent, on conclut selon que
l'égalité obtenue est fausse ou toujours vraie, et on interprète avec les droites.
\etapeexemples
Résoudre chaque système, puis interpréter le résultat avec les droites associées.
\begin{center}
  a)~${\sys{2x-y &= 3 && (L_{1})\\ -4x+2y &= 1 && (L_{2})}}$ \qquad
  b)~${\sys{x+2y &= 4 && (L_{1})\\ 3x+6y &= 12 && (L_{2})}}$
\end{center}

\etape{À vous de jouer}
Résoudre chaque système, puis interpréter le résultat avec les droites associées.
\begin{center}
  a)~${\sys{3x+y &= 2 && (L_{1})\\ -6x-2y &= 5 && (L_{2})}}$ \qquad
  b)~${\sys{x-3y &= 2 && (L_{1})\\ -2x+6y &= -4 && (L_{2})}}$
\end{center}

\end{methodebox}

\afaire{exercices 19 à 21}

\end{document}
